\documentclass[letterpaper,12pt]{article}
\usepackage{amsmath, amssymb, geometry, titlesec, esint, listings, xcolor, graphicx, float, caption}

\geometry{margin=1in}

%Make section titles smaller
\titleformat{\section}
    {\large\bfseries}
    {\thesection.}{0.5em}{}

\titleformat{\subsection}
    {\normalsize\bfseries}
    {\thesubsection.}{0.5em}{}

%Customize vertical spacing
\titlespacing*{\section}{0pt}{6pt}{3pt}
\titlespacing*{\subsection}{0pt}{6pt}{3pt}

\setlength{\textfloatsep}{5pt}
\setlength{\floatsep}{5pt}
\setlength{\intextsep}{5pt}

%\linespread{0.97}

\setlength{\parindent}{0pt} % no indentation

%Figure caption sizes
\captionsetup[figure]{justification=raggedright, font=small, labelfont=bf, skip=0pt}

\lstset{
    language=Python,
    basicstyle=\ttfamily\small,
    keywordstyle=\color{blue!60!black},
    commentstyle=\color{green!50!black},
    stringstyle=\color{teal!60!black},
    numbers=left,
    numberstyle=\tiny\color{gray},
    stepnumber=1,
    breaklines=true,
    frame=single,
    tabsize=4,
    showstringspaces=false
}

\makeatletter
\g@addto@macro\normalsize{%
  \setlength\abovedisplayskip{6pt}%
  \setlength\belowdisplayskip{6pt}%
  \setlength\abovedisplayshortskip{6pt}%
  \setlength\belowdisplayshortskip{6pt}%
}
\makeatother

\begin{document}

%---------------------------------------------------------------------------------------------------------------

%INFO

Christian DiPietrantonio \\
ME 5311: Computational Methods to Viscous Flows \\
\today

\vspace{6pt}

%---------------------------------------------------------------------------------------------------------------

%TITLE
{\large \textbf{Computer Assignment 02:} \normalsize Obtaining a Numerical Solution to the Zero Pressure Gradient Laminar Boundary Layer Equations}

%---------------------------------------------------------------------------------------------------------------

%DESCRIPTION OF THE PROBLEM
\section{Description of the Problem}

TBD

%---------------------------------------------------------------------------------------------------------------

%DESCRIPTION OF THE NUMERICAL METHOD
\section{Description of the Numerical Method}

Simplifications:

\begin{equation}
\boxed{
    Re_L = \frac{U_{\infty} L}{\nu_{\infty}}
}
\end{equation}

\begin{equation}
\boxed{
    C = \frac{1}{Re_L} \left( \frac{L}{\delta} \right)^2
}
\end{equation}

Non-Dimensional system:

\begin{equation}
\boxed{
    \frac{\partial u^*}{\partial x^*} + \frac{\partial v^*}{\partial y^*} = 0
}
\end{equation}

\begin{equation}
\boxed{
    u^* \frac{\partial u^*}{\partial x^*} + v^* \frac{\partial u^*}{\partial y^*} = C \frac{\partial}{\partial y^*} \left(\nu^* \frac{\partial u^*}{\partial y^*} \right)
}
\end{equation}

Transformation from $y$ to $\eta$:
\begin{equation}
\boxed{
    J = \frac{\partial y}{\partial \eta} = \frac{y_{max} s \cosh(s \eta)}{\sinh(s)}
}
\end{equation}

\begin{equation}
\boxed{
    J_{\eta} = \frac{\partial J}{\partial \eta} = \frac{s^2 y_{max} \sinh(s \eta)}{\sinh(s)}
}
\end{equation}

\begin{equation}
\boxed{
    \frac{\partial}{\partial y} = \frac{1}{J} \frac{\partial}{\partial \eta}
}
\end{equation}

\begin{equation}
\boxed{
    \frac{\partial^2}{\partial y^2} = \frac{1}{J^2} \frac{\partial^2}{\partial \eta^2} - \frac{J_{\eta}}{J^3} \frac{\partial}{\partial \eta}
}
\end{equation}


\begin{equation}
\boxed{
    J^* = \frac{\partial y^*}{\partial \eta} = \frac{J}{\delta} = \frac{y_{max} s \cosh(s \eta)}{\delta \sinh(s)}
}
\end{equation}

\begin{equation}
\boxed{
    J_{\eta}^* = \frac{\partial J^*}{\partial \eta} = \frac{s^2 y_{max} \sinh(s \eta)}{\delta \sinh(s)}
}
\end{equation}

\begin{equation}
\boxed{
    \frac{\partial}{\partial y^*} = \frac{1}{J^*} \frac{\partial}{\partial \eta}
}
\end{equation}

\begin{equation}
\boxed{
 \frac{\partial^2}{\partial {y^*}^2} = \frac{1}{{J^*}^2} \frac{\partial^2}{\partial \eta^2} - \frac{J_{\eta}^*}{{J^*}^3} \frac{\partial}{\partial \eta}
}
\end{equation}

Final system:

\begin{equation}
\boxed{
    \frac{\partial v^*}{\partial \eta} = -J^* \frac{\partial u^*}{\partial x^*}
}
\end{equation}

\begin{equation}
\boxed{
    u^* \frac{\partial u^*}{\partial x^*} + A \frac{\partial u^*}{\partial \eta} + B \frac{\partial^2 u^*}{\partial \eta^2} = 0
}
\end{equation}

\begin{equation}
\boxed{
    A = \frac{v^*}{J^*} - \frac{C}{{J^*}^2} \frac{\partial \nu^*}{\partial \eta} + \frac{C \nu^* J_{\eta}^*}{{J^*}^3}
}
\end{equation}

\begin{equation}
\boxed{
    B = - \frac{C \nu^*}{{J^*}^2}
}
\end{equation}




%---------------------------------------------------------------------------------------------------------------

%PRESENTATION OF RESULTS
\section{Presentation of Results}

TBD

%---------------------------------------------------------------------------------------------------------------

%DISCUSSION OF RESULTS
\section{Discussion of Results}

\subsection{General description}

TBD

\subsection{Accuracy and stability}

TBD

%---------------------------------------------------------------------------------------------------------------

%APPENDIX A - copy of program listing
\newpage
\appendix
\section{Copy of Program Listing}
\label{appendix-code}

% \lstinputlisting[
%     language=Python
% ]{<insert python code>}

%---------------------------------------------------------------------------------------------------------------

\end{document}
